% THETA3-B22 contribution by /root/definite_bridge, 21 November 2026.
\subsection{THETA3-B22: explicit theta-null and tangent normalization bounds}\label{B22}

The scalar in B19 need not be left with an unknown absolute
CM-period size.  An explicitly chosen even reference and an
explicitly nonzero odd first jet give an algebraic normalization
whose analytic factor has a direct lower certificate.  This
does not bound the heights or denominators of the resulting
algebraic coordinates, and does not assert that they are an
integral saturated frame at a specified prime.

\paragraph{The two references and their Gaussian tails.}
Fix $Z=X+iY$ in the Siegel upper half space, a dimension $g$,
and a nonzero specified tangent vector $u\in\C^g$.
Fix a certified number $\lambda>0$ with $Y\geq\lambda I_g$.
At a rational CM point the matrix entries are algebraic, so
one can choose a positive rational $\lambda$ and certify this
inequality by the positive principal minors of $Y-\lambda I_g$.
Choose an index $j$ with $u_j\ne0$ and put $q=e_j$.
Let $N>2$ be an even multiple of a level containing all the
previously prescribed finite inputs.  The two functions
\[
 \varphi_{\rm ev}=1_{N\Z^g},\qquad
 \varphi_{\rm odd}=1_{q+N\Z^g}-1_{-q+N\Z^g}
 \tag{B22a}
\]
belong to $W_N=\operatorname{Fun}(N^{-1}\Z^g/N\Z^g,\Q)$.
Their two odd cosets are distinct.  Use the standard theta
sum and its completed version from (B19b).  At zero the
completion is one, and its first derivative is zero.

For the even theta-null $E_N$ one has
\[
 E_N=\sum_{m\in\Z^g}e^{\pi iN^2m^tZm},\qquad
 |E_N-1|\leq
       \left(\frac{1+r}{1-r}\right)^g-1,
 \quad r=e^{-\pi\lambda N^2}.
 \tag{B22b}
\]
Indeed $m^2\geq |m|$ for an integer, so
$\sum_{m\in\Z}r^{m^2}\leq(1+r)/(1-r)$; taking the
product over the coordinates and deleting zero proves the
bound.  We may increase $N$ until its right side is at most
$1/2$, obtaining
\[
                         \frac12\leq |E_N|\leq\frac32.
 \tag{B22c}
\]

The odd first derivative has the exact expansion
\[
 H_N=D_u\widehat\vartheta_{\varphi_{\rm odd}}(Z,0)
       =4\pi i\sum_{m\in\Z^g}
          ((q+Nm)^tu)e^{\pi i(q+Nm)^tZ(q+Nm)}.
 \tag{B22d}
\]
To see the factor four and the sign, differentiate both
theta series and replace $m$ by $-m$ in the second.  The
linear factor changes sign and the quadratic factor does
not, so the two differentiated series add.
The zero term is
\[
 a_0=4\pi i(q^tu)e^{\pi iq^tZq},\qquad
 |a_0|=4\pi|q^tu|e^{-\pi q^tYq}>0.
\]
Put $a=\pi\lambda N^2/4$ and $t=e^{-a}$.
Since $N\geq2\|q\|$ and $\|m\|\geq1$ for $m\ne0$,
\[
 \|q+Nm\|\geq\frac N2\|m\|,\qquad
 \|q+Nm\|\leq\frac{3N}{2}\|m\|.
\]
Therefore
\[
 |H_N-a_0|\leq
  12\pi gN\|u\|\,
       \frac{t}{(1-t)^2}
       \left(\frac{1+t}{1-t}\right)^{g-1}
       =:T_N.
 \tag{B22e}
\]
For completeness, the initial bound is
$6\pi N\|u\|\sum_{m\ne0}\|m\|e^{-a\|m\|^2}$.
Use $\|m\|\leq\sum_k|m_k|$, then separate the sums.
In the distinguished coordinate,
\[
 \sum_{b\in\Z}|b|t^{b^2}
       \leq2\sum_{b\geq1}bt^b
       =\frac{2t}{(1-t)^2};
\]
each of the other coordinates contributes at most
$(1+t)/(1-t)$.  Summing the $g$ possibilities gives (B22e).

Increase $N$ through even multiples of the prescribed level
until, in addition to (B22c),
\[
                 T_N\leq
                 2\pi|q^tu|e^{-\pi q^tYq}.
 \tag{B22f}
\]
This process terminates: the right side is a fixed positive
number, whereas $T_N$ is a polynomial factor in $N$ times a
Gaussian tending to zero.  It is a finite certification
procedure with explicit input data.  The result is
\[
 |H_N|\geq B(Z,u,q):=
               2\pi|q^tu|e^{-\pi q^tYq}>0.
 \tag{B22g}
\]
No theta nonvanishing theorem or CM-period estimate is used
in this calculation.

\paragraph{The corresponding algebraic theta and tangent frames.}
Now use the actual CM torus of B19 and its algebraic finite
theta isomorphism $J_N$.
First normalize it algebraically so that
\[
                         J_N\varphi_{\rm ev}(0)=1.
 \tag{B22h}
\]
This is possible because its value is algebraic and nonzero:
the analytic value is $E_N\ne0$, and the two theta maps differ
by a nonzero scalar.  Dividing $J_N$ by its previous nonzero
value is an allowed algebraic change of frame.  With the
standard oscillator normalization, the new comparison scalar
is exactly $c_N=E_N$.  For a prescribed overall oscillator
scalar $\upsilon$, it is $c_N=\upsilon E_N$ instead; that
known scalar is retained.

Fix the rational $E$-linear CM decomposition from (B19a).
Choose one nonzero Betti eigenline direction in its first
$B(d)$ copy and let $u=u^w$ be its image in the $w$
coordinate of this actual torus.  The same coordinate vector
in the original theta variable is denoted $u^z$.
The selected complex tangent line is the complexification of
an algebraic CM eigenspace.  Its first-jet functional
\[
       \ell(v)=D_v(J_N\varphi_{\rm odd})(0)
\]
is algebraic on that eigenspace, because the section, line,
and symmetric origin first-jet frame are algebraic.
The exact substitution $z=Nw$ gives
\[
                          c_N\ell(u^w)=N\upsilon H_N.
 \tag{B22i}
\]
The right side is nonzero by (B22g).  Choose an algebraic
generator $v$ of the line with $\ell(v)=1$, by dividing any
algebraic generator by its nonzero algebraic $\ell$-value.
Writing $u^w=p_Nv$ now gives the exact identity
\[
                    \frac{c_Np_N}{N}=\upsilon H_N.
 \tag{B22j}
\]
Thus the tangent period of B19, in this specified algebraic
frame, has not been assumed bounded: its product with the
theta-null normalization has been identified with a directly
bounded analytic first derivative.

Use the same generator of the CM factor in its second
repeated copy, transported by the fixed algebraic isogeny.
The corresponding Betti directions are obtained from the
same Betti generator in the two copies.  They have the same
comparison scalar $p_N$.  Any original pair of selected
ball directions is obtained from these two by an invertible
matrix over $\overline\Q$: both are bases in the same
two-dimensional Betti eigenspace, with the rational
$E$-decomposition and the rational negative-line coordinates
fixed.  This is precisely the matrix $M_i$ in (B19d).
Consequently (B22j) can be imposed for each input type while
preserving B19's common scalar for both repeated directions.
The reference depends on the torus and the chosen direction,
not on the source character or the finite original input.

\paragraph{The resulting period certificate and its arithmetic cost.}
Apply this normalization to the two factors in (B19j).
The common constant becomes exactly
\[
                  C_x=\kappa_x\upsilon_1\upsilon_2
                                      H_{N_1}H_{N_2}.
\]
It follows that
\[
 |C_x|^2\geq
   |\kappa_x\upsilon_1\upsilon_2|^2
       \prod_{i=1}^2
        \left(2\pi|q_i^tu_i|
                      e^{-\pi q_i^tY_iq_i}\right)^2.
 \tag{B22k}
\]
Here $\kappa_x$ is in the fully specified frame of B19;
under the algebraic Betti choices of B21 it is computed by
the indicated nonzero determinant and square-root factors,
with the common universal oscillator factors retained.
Arbitrary original complex rescalings cannot be omitted.

For a nonzero algebraic finite sum $b$ in (B19j), suppose
$b\in K$, $[K:\Q]=d$, $Db$ is integral, and the other
$d-2$ conjugates are bounded by $M>0$ in B19's nonreal
embedding convention.  Combining (B19k) and (B22k) gives
\[
\begin{split}
 |\mathcal B_e(f)|^2\ \geq\ &
 |v_Hv_1^Tv_2^T\kappa_x\upsilon_1\upsilon_2|^2
       D^{-d}M^{-(d-2)}\\
 &{}\times\prod_{i=1}^2
        \left(2\pi|q_i^tu_i|
                      e^{-\pi q_i^tY_iq_i}\right)^2.
\end{split}
 \tag{B22l}
\]
For an orthonormal test vector use the $v_H^{1/2}$ factor
specified in B19 instead.  In combination with Q11 this
also gives a lower certificate for the geometric canonical
norm when its stated chart and evaluation data are supplied.

The constants $D,M$ refer to the algebraic array in these
new frames.  Dividing by the algebraic reference value in
(B22h), and by $\ell(v_0)$ when choosing $v$, can introduce
denominators and large conjugates.  Those factors must be
included in $D,M$ and in any prime-by-prime valuation bound.
The normalized tangent frame need not be integral or
saturated, so B20 or the archive's local divisibility
theorems require their integral hypotheses to be checked
again or their explicit change-of-frame costs to be applied.
No uniform arithmetic height bound follows from the analytic
Gaussian estimate.  The improvement is the removal of an
unspecified absolute CM-period factor from the normalized
finite-period certificate, with its arithmetic costs assigned
to actual algebraic data.\hfill$\square$
